% ============================================================================
\section{Hasil}
\label{sec:hasil}
% ============================================================================

\catatan{Belum ada hasil. Bagian ini hanya kerangka urutan pelaporan, tabel kosong, dan gambar pengganti. Isi setelah E0 lulus dan E1--E4 dijalankan sesuai pra-registrasi. Hipotesis yang terbantah dilaporkan sebagai terbantah, dan analisis di luar pra-registrasi diberi label eksploratori.}

Urutan pelaporan mengikuti Tabel~\ref{tab:lacak}. Validasi instrumen dilaporkan lebih dulu. Setiap RQ dibuka dengan metrik utama dan putusan atas hipotesisnya, lalu metrik diagnostik.

\subsection{Validasi instrumen}
\label{subsec:hasil-e0}
\isian{Selisih $\K^{\emas}$ terhadap GT1--GT5 menurut hasil adjudikasi; matriks acuan; Tabel~\ref{tab:cakupan} yang sudah terisi; \en{mutation score} global dan terfokus-$\Delta$; Krippendorff's $\alpha$ untuk $\Delta^{\emas}$; jumlah koreksi emas.}

\subsection{Adaptivitas representasi (RQ1)}
\label{subsec:hasil-rq1}
\isian{RA per kelompok dengan CI 95\% dan putusan H1; perubahan dan kebijakan yang gagal beserta alasannya; riwayat \en{engine} yang dipetakan ke regulasi pemicunya; akurasi resolusi konflik.}

\begin{table}[t]
\centering
\caption{Adaptivitas representasional (E1)}
\label{tab:hasil-rq1}
\scriptsize
\begin{tabularx}{\columnwidth}{@{}L P{0.5cm} P{1.0cm} P{1.9cm}@{}}
\toprule
\textbf{Kelompok} & $n$ & \textbf{Berhasil} & \textbf{RA (CI 95\%)} \\
\midrule
Perubahan nyata 2023--2026 & \isian{} & \isian{} & \isian{} \\
Perubahan tak terlihat (H1) & \isian{} & \isian{} & \isian{} \\
Kebijakan perusahaan nyata & \isian{} & \isian{} & \isian{} \\
Kebijakan buatan & \isian{} & \isian{} & \isian{} \\
\bottomrule
\end{tabularx}
\end{table}

\subsection{Akuisisi satu langkah (RQ2)}
\label{subsec:hasil-rq2}
\isian{SER, CAR, dan ABS per kondisi, per lapisan, dan per jenis perubahan dengan CI; selisih CAR dan non-inferioritas SER terhadap B3 dan B4 beserta putusan H2a dan H2b; penilaian relevansi; strata kontrafaktual; \en{pass@1}, \en{1@m}, dan \en{pass@m}; metrik diagnostik.}

\begin{table}[t]
\centering
\caption{Keamanan dan kegunaan akuisisi per kondisi (E2)}
\label{tab:hasil-rq2}
\scriptsize
\begin{tabularx}{\columnwidth}{@{}L P{1.45cm} P{1.45cm} P{1.1cm}@{}}
\toprule
\textbf{Kondisi} & \textbf{SER (CI)} & \textbf{CAR (CI)} & \textbf{ABS} \\
\midrule
Usulan, lapisan regulasi & \isian{} & \isian{} & \isian{} \\
Usulan, lapisan perusahaan & \isian{} & \isian{} & \isian{} \\
B3 (\en{pass@1}, G1--G2) & \isian{} & \isian{} & \isian{} \\
B4 (suara terbanyak) & \isian{} & \isian{} & \isian{} \\
AB2--AB5 & \isian{} & \isian{} & \isian{} \\
\bottomrule
\end{tabularx}
\end{table}

\begin{figure}[t]
\centering
\gambar{Setiap kondisi digambar sebagai titik dengan CI. Sumbu-x adalah cakupan, yaitu CAR $+$ SER. Sumbu-y adalah risiko selektif, yaitu SER$/$(CAR $+$ SER), atau laju kesalahan di antara pembaruan yang diterapkan.}
\caption{Keamanan dan kegunaan akuisisi per kondisi.}
\label{fig:rq2}
\end{figure}

\subsection{Adaptivitas operasional berantai (RQ3)}
\label{subsec:hasil-rq3}
\isian{OAR dan putusan H3a; kesalahan bawaan, baru, dan terpicu per peristiwa beserta putusan H3b; rasio perambatan; jumlah peristiwa yang diterapkan tanpa suntingan manusia dan waktunya; \en{drift} akhir dalam rupiah; anomali \en{knowledge base}; perbandingan dengan B1.}

\begin{figure*}[t]
\centering
\gambarlebar{Satu batang bertumpuk per peristiwa, diurutkan kronologis di sumbu-x dan diberi label jenis perubahan. Batang menunjukkan kesalahan bawaan, baru, dan terpicu. Penanda menunjukkan status A2, A3$_t$, dan A4, dan garis menunjukkan rentang antarpengulangan.}
\caption{Hasil \en{regulation replay}.}
\label{fig:replay}
\end{figure*}

\subsection{Analisis dampak lintas lapis (RQ4)}
\label{subsec:hasil-rq4}
\isian{\en{Recall} dan $F_2$ dengan CI pada kedua masukan ($\Delta^{\emas}$ dan $\hat{\Delta}$) beserta putusan H4a dan H4b; perbandingan dengan B5, termasuk garis dasar \emph{tandai semua}; Cohen's $\kappa$; satu contoh penelusuran dari graf.}

\subsection{Ablasi}
\label{subsec:hasil-ablasi}
\isian{Perubahan SER dan CAR untuk setiap ablasi. AB5 menjawab dengan angka mengapa \en{knowledge graph} diperlukan.}
